\documentclass{article}

% ── Packages ──────────────────────────────────────────────
\usepackage[margin=1in]{geometry}
\usepackage{times}
\usepackage[T1]{fontenc}
\usepackage[utf8]{inputenc}
\usepackage{amsmath, amssymb, mathtools}
\usepackage{booktabs}
\usepackage{graphicx}
\usepackage{caption}
\usepackage{float}
\usepackage{listings}
\usepackage{xcolor}
\usepackage{microtype}
\usepackage{hyperref}
\usepackage{natbib}

\bibliographystyle{plainnat}

\lstset{
    basicstyle=\small\ttfamily,
    backgroundcolor=\color{gray!10},
    frame=single,
    breaklines=true,
    captionpos=b,
}

\hypersetup{
    colorlinks=true,
    linkcolor=blue,
    citecolor=blue,
    urlcolor=blue,
}

% ══════════════════════════════════════════════════════════

\title{\textbf{Recursive Convergent Inference (RCI):}\\
Bottom-Up Module Expansion via Output Convergence}

\author{
    Alex Anokhin \\
    \textit{Independent Researcher} \\
    \href{mailto:olanokhin@gmail.com}{olanokhin@gmail.com} \\[4pt]
    \small Code \& Colab:
    \href{https://github.com/olanokhin/rci-inference}
         {github.com/olanokhin/rci-inference}
}

\date{March 2026}

% ══════════════════════════════════════════════════════════
\begin{document}
\maketitle

% ── ABSTRACT ──────────────────────────────────────────────
\begin{abstract}
We propose \textbf{Recursive Convergent Inference (RCI)}, an
architectural principle for neural network inference in which the
set of active computational modules expands monotonically from a
minimal seed subset until empirical convergence of the model's
next-token output distribution.
Unlike existing adaptive computation methods that determine
\textit{when to halt} a fixed computation, RCI determines
\textit{when additional computation is warranted}---progressively
incorporating additional modules based on inter-module affinities
until output stability is reached.
Inference terminates when a hybrid distributional stability metric
falls below a threshold, with additional safeguards against
premature convergence.
Crucially, RCI requires no learned halting signal, external
verifier, or task-complexity pre-estimator; stopping emerges
directly from the internal stability of the model's own output
distribution.
As each inference process is self-contained, RCI naturally enables
parallel multi-task execution over shared read-only weights.
This offers a complementary perspective to existing test-time
compute scaling by shifting from \textit{external scaling}
(longer outputs, multiple samples, verifier-guided search) to
\textit{internal scaling} over the active parameter subgraph.

We instantiate RCI on OLMoE-1B-7B-0924---a fully open model with
64 discrete experts.
Empirical evaluation on reasoning benchmarks of varying difficulty
(GSM8K, MATH, MMLU hard subsets, $n{=}150$) demonstrates that all
pairwise difficulty comparisons are statistically significant:
hard tasks accumulate significantly greater total KL divergence
than easy (Mann-Whitney $U{=}1677$, $p{=}0.002$) and medium tasks
($U{=}2106$, $p{<}0.001$), with a statistically significant
monotonic correlation (Spearman $\rho{=}0.22$, $p{=}0.007$).
Notably, RCI's complexity metric diverges from human-defined
difficulty labels, suggesting it captures \textit{model-relative
computational demand} rather than task difficulty in the abstract.
We release a fully reproducible implementation as a Google Colab
notebook alongside this paper.

\textbf{Keywords:} adaptive inference, mixture-of-experts,
output convergence, dynamic subgraph expansion, test-time compute,
efficient inference
\end{abstract}

% ── 1. INTRODUCTION ───────────────────────────────────────
\section{Introduction}
\label{sec:intro}

The computational cost of inference in large neural networks is
conventionally fixed by architecture: a model with $L$ layers and
$N$ parameters applies largely the same amount and structure of
computation to every input, regardless of whether the task
requires synthesizing multi-step reasoning chains or answering a
simple factual query.
This rigidity constitutes a fundamental inefficiency---a direct
retrieval question consumes the same resources as a complex
mathematical proof.

Recent advances have pursued two complementary directions.
\textit{External test-time scaling} allocates additional compute
by extending token sequences or sampling multiple trajectories,
selecting among them via process reward models or majority
voting~\citep{snell2024scaling, muennighoff2025s1}.
\textit{Adaptive computation} methods reduce wasted compute by
learning early exits or layer skips from a fixed graph, as in
Adaptive Computation Time (ACT)~\citep{graves2016act},
Universal Transformers~\citep{dehghani2019universal},
PonderNet~\citep{banino2021pondernet}, and
Mixture-of-Depths~\citep{raposo2024mod}.
Both approaches operate \textbf{top-down}: they begin with a
maximal computational graph and learn when to stop or what to
remove.
The symmetric direction---beginning from a minimal active
parameter subset and expanding it bottom-up as a function of
task complexity---has remained comparatively underexplored.

We propose \textbf{Recursive Convergent Inference (RCI)}, an
inference principle that addresses this gap.
Rather than fixing the active parameter set and learning when to
halt, RCI begins from a minimal seed subset of modules and expands
it monotonically ($M_n \subseteq M_{n+1}$) until the model's
next-token output distribution converges.
The stopping criterion is intrinsic: additional modules no longer
meaningfully alter the output distribution, without relying on an
explicit learned halting signal or external evaluator.

This formulation induces several properties.
First, computation is \textit{task-adaptive by construction}:
simple inputs converge early over small subgraphs, while complex
reasoning recruits progressively larger portions of the network.
Second, each inference process is \textit{self-contained},
operating over shared read-only weights and potentially enabling
parallel multi-task execution.
Third, it is \textit{estimator-free}: termination arises from the
distributional stability of the model's own output.
More broadly, this suggests that computation in large networks can
be allocated dynamically at the level of parameter activation,
rather than being fixed by architecture or governed by external
sampling strategies.

\paragraph{Contributions.} We make the following contributions:
\begin{enumerate}
    \item \textbf{Recursive Convergent Inference (RCI)}---a novel
    inference principle formulating computation as monotonic
    module-set expansion ($M_n \subseteq M_{n+1}$) until empirical
    output convergence, providing a bottom-up counterpart to
    existing top-down adaptive methods.

    \item \textbf{Module affinity graph}---a lightweight static
    graph derived from expert co-activation patterns in MoE models,
    used to guide structured expansion of the active module set.

    \item \textbf{Convergence stopping criterion}---a hybrid output
    stability metric combining KL divergence, cosine similarity,
    and output entropy over a rolling window, conditioned on a
    Confidence Margin to mitigate false convergence without
    external supervision.

    \item \textbf{Empirical validation}---an open, reproducible
    implementation on OLMoE-1B-7B-0924 demonstrating that the
    computational effort required for convergence (measured via
    cumulative KL divergence across expansion steps) scales with
    task complexity.
    All pairwise difficulty comparisons are statistically
    significant across $n{=}150$ examples (50 per tier---GSM8K,
    MATH algebra, MMLU hard subsets): hard tasks accumulate
    significantly greater cumulative KL divergence than easy
    tasks (Mann-Whitney $U{=}1677$, $p{=}0.002$) and medium
    tasks ($U{=}2106$, $p{<}0.001$), with a statistically
    significant monotonic correlation (Spearman $\rho{=}0.22$,
    $p{=}0.007$, $n{=}150$).
    Notably, RCI's complexity metric diverges from human-defined
    difficulty labels, indicating model-relative computational
    demand rather than task difficulty in the abstract.
\end{enumerate}

% ── 2. RELATED WORK ───────────────────────────────────────
\section{Related Work}
\label{sec:related}

RCI lies at the intersection of five active research directions:
test-time compute scaling, adaptive computation, sparse
mixture-of-experts routing, fixed-point iterative inference, and
graph-based computation.

\paragraph{Test-time compute scaling.}
Scaling computation at inference time has emerged as a primary
axis for improving reasoning in large language
models~\citep{snell2024scaling, agarwal2025art}.
Existing approaches fall into two categories.
\textit{Serial scaling} increases compute by extending token
sequences (Chain-of-Thought, iterative revision).
\textit{Parallel scaling} samples multiple output trajectories
and selects among them via verifiers or reward models
(Best-of-$N$, Tree Search, MCTS-based planning).
Self-consistency~\citep{wang2022selfconsistency} samples multiple
independent outputs and selects by majority vote.
Simple test-time scaling~\citep{muennighoff2025s1} demonstrates
that SFT on a small set of reasoning traces combined with budget
forcing achieves strong scaling behaviour.
RCI differs by scaling computation \textit{internally}---expanding
the active parameter subgraph per forward pass without altering
context length or requiring any external verifier.

\paragraph{Adaptive computation and halting.}
ACT~\citep{graves2016act} introduces a learned scalar halting
probability for recurrent networks.
The Universal Transformer~\citep{dehghani2019universal} extends
this to transformers via weight-shared recurrent layers.
PonderNet~\citep{banino2021pondernet} reformulates halting
probabilistically.
Mixture-of-Depths~\citep{raposo2024mod} learns per-token binary
skip decisions.
More recent work includes AdaPonderLM~\citep{adaponderlm2025} and
ADEPT~\citep{adept2026}.
Answer Convergence~\citep{liu2025answerconv} detects when outputs
stabilise across reasoning steps---the closest prior use of a
convergence signal for stopping.
All of the above focus on \textit{halting}; RCI asks the symmetric
question of \textit{when additional computation is warranted}.

\paragraph{Mixture-of-Experts.}
The Switch Transformer~\citep{fedus2021switch} established sparse
MoE as a practical scaling paradigm.
OLMoE~\citep{muennighoff2024olmoe} demonstrated that open MoE
models with 64 experts and top-8 routing match larger dense
counterparts.
OpenMoE~\citep{xue2024openmoe} confirmed strong expert
specialisation during pre-training.
Standard and dynamic MoE methods make a single routing decision
per forward pass.
RCI treats the active expert set as a dynamically expanding
frontier across recursive steps, repurposing co-activation
statistics as a static affinity graph---a structural use of MoE
routing information not present in prior work.

\paragraph{Fixed-point and iterative inference.}
Deep Equilibrium Models~\citep{bai2019deq} formulate inference
as finding a fixed point of a weight-tied layer.
Looped Transformers and recurrent-depth architectures demonstrate
that iterating a single parameterised block can scale test-time
computation.
SELF-Transformer~\citep{mathur2025self} applies fixed-point
iteration within a single self-attention layer.
Inner Thinking Transformer~\citep{chen2025inner} interprets
depth as internal reasoning steps.
All existing fixed-point methods refine representations within a
\textit{fixed} computational structure; RCI differs in that
convergence is achieved through \textit{expansion} of the active
module set.

\paragraph{Hardware-efficient sparse operations.}
Flash-KMeans~\citep{flashkmeans2025} demonstrates that IO-aware
kernel engineering can yield 30--200$\times$ speedups for sparse
operations previously considered hardware-unfriendly, suggesting
a concrete engineering path toward reducing RCI's wall-clock
overhead without algorithmic changes.

\paragraph{Summary of distinctions.}
Table~\ref{tab:related} summarises the key distinctions.

\begin{table}[ht]
\centering
\caption{Comparison of RCI with related methods.}
\label{tab:related}
\small
\begin{tabular}{lccccc}
\toprule
Method & Dyn.\ Depth & Monotonic Width & Conv.\ Stop &
No Pre-est. & Parallel \\
\midrule
ACT / UT / PonderNet & \checkmark & $\times$ &
$\times$ & $\times$ & $\times$ \\
Mixture-of-Depths & \checkmark & $\times$ &
$\times$ & $\times$ & $\times$ \\
Answer Convergence & \checkmark & $\times$ &
\checkmark & $\times$ & $\times$ \\
Switch / OLMoE & $\times$ & $\times$ &
$\times$ & $\times$ & $\times$ \\
SELF-Transformer & \checkmark & $\times$ &
\checkmark & \checkmark & $\times$ \\
Best-of-$N$ / MCTS & $\times$ & $\times$ &
ext.\ budget & $\times$ & $\times$ \\
\midrule
\textbf{RCI (ours)} & \checkmark & \checkmark &
\checkmark & \checkmark & \checkmark \\
\bottomrule
\end{tabular}
\end{table}

% ── 3. METHOD ─────────────────────────────────────────────
\section{Recursive Convergent Inference}
\label{sec:method}

\subsection{Core Principle}

RCI formulates inference as monotonic module-set expansion until
empirical output convergence.
In contrast to existing top-down adaptive methods, RCI operates
bottom-up: computation begins at a minimal seed subset and grows
as warranted by the task.

\subsection{Seed Initialisation}

Seeds define the initial subset of active computational modules.
In the PoC, seeds are initialised as the top-8 most frequently
activated experts on the first standard forward pass.
In the full architecture, a lightweight input classifier
(embedding + linear router) maps task type to domain-specific
module clusters, trained with an auxiliary loss including
$\lambda_4 \cdot \mathrm{Var}(M_{\mathrm{active}})$ to encourage
distinct subgraph activation across task types.

\subsection{Recursive Expansion}

The affinity graph encodes co-activation frequency between modules:

\[
\mathrm{affinity}[i][j] \mathrel{+}= 1
\;\text{ if experts } i,j
\text{ are co-activated within the same sequence.}
\]

The expansion rule is:
\begin{align}
M_0 &= \text{seeds} \\
M_{n+1} &= M_n \cup
\mathrm{top\text{-}k}\!\left(\mathrm{neighbors}(M_n),\,
\mathrm{weighted\ by\ affinity}\right),
\quad M_n \subseteq M_{n+1}
\end{align}

Implementation proceeds via a gating mask:
\begin{lstlisting}[language=Python]
gating_logits[experts_not_in_M] = -inf
current_k = min(base_k * (2 ** step), total_experts)
\end{lstlisting}

\subsection{Convergence Criterion}

\begin{equation}
\Delta_n = \alpha \cdot \mathrm{KL}(P_n \,\|\, P_{n-1})
+ \beta \cdot \cos(h_n, h_{n-1})
+ \gamma \cdot H(P_n)
\end{equation}

\begin{equation}
\text{Convergence: } \frac{1}{W}\sum_{i=n-W+1}^{n}\Delta_i
< \varepsilon \;\;\text{and}\;\;
\bigl[P(\text{top-1}) - P(\text{top-2})\bigr] \geq \theta
\end{equation}

The Confidence Margin guards against false convergence at local
plateaus without requiring an external estimator.

\subsection{Training Objective}

\begin{equation}
\mathcal{L} = \mathcal{L}_{\mathrm{task}}
+ \lambda_1 \mathcal{L}_{\mathrm{conv}}
+ \lambda_2 \mathcal{L}_{\mathrm{sparse}}
+ \lambda_3 \,\mathrm{MI}(\mathrm{seeds},\,M_{\mathrm{final}})
+ \lambda_4 \,\mathrm{Var}(M_{\mathrm{active}})
\label{eq:loss}
\end{equation}

% ── 4. EXPERIMENTS ────────────────────────────────────────
\section{Experiments}
\label{sec:experiments}

\subsection{Setup}

\paragraph{Model.}
OLMoE-1B-7B-0924~\citep{muennighoff2024olmoe}: 64 experts per
MoE layer, top-8 routing, 1B active / 7B total parameters,
Apache 2.0 licence.
All experiments on a single NVIDIA T4 GPU (16 GB) via Google
Colab.

\paragraph{Affinity graph.}
500 calibration examples from C4~\citep{raffel2020c4}.
Co-activation counts are L$_\infty$-normalised to $[0,1]$.
Construction requires approximately 10 minutes on T4.

\paragraph{RCI loop.}
Seeds: top-8 most activated experts on first pass.
Expansion: BFS over affinity graph, $k{=}4$ neighbours per step.
Gating: dynamic top-$k$ doubling ($8{\to}16{\to}32{\to}64$).
Convergence: $\varepsilon{=}0.093$ (calibrated as the 20th
percentile of KL values over 5 calibration tasks),
$W{=}3$, max steps $= 15$.

\paragraph{Benchmarks.}
\begin{itemize}
    \item \textbf{Easy}: GSM8K~\citep{cobbe2021gsm8k} ---
    grade school mathematics word problems
    \item \textbf{Medium}: MATH algebra~\citep{hendrycks2021math}
    \item \textbf{Hard}: MMLU hard subsets~\citep{hendrycks2021mmlu}
    --- college physics, chemistry, abstract algebra, formal logic
\end{itemize}
50 examples per tier, $n{=}150$ total.

\subsection{Results}

\begin{table}[ht]
\centering
\caption{Cumulative KL divergence (AUC) by difficulty tier
($n{=}150$, 50 per tier).}
\label{tab:results}
\begin{tabular}{llccc}
\toprule
Difficulty & Benchmark & $n$ & Mean AUC & Std \\
\midrule
Easy & GSM8K & 50 & 10.728 & 2.808 \\
Medium & MATH (algebra) & 50 & 8.956 & 1.688 \\
Hard & MMLU hard & 50 & 11.987 & 2.537 \\
\bottomrule
\end{tabular}
\end{table}

All pairwise Mann-Whitney U tests are statistically significant:
hard vs.\ easy ($U{=}1677$, $p{=}0.002$),
hard vs.\ medium ($U{=}2106$, $p{<}0.001$),
easy vs.\ medium ($U{=}1788$, $p{<}0.001$).
The overall Spearman rank correlation is
$\rho{=}0.22$ ($p{=}0.007$, $n{=}150$).

\begin{figure}[h]
\centering
\includegraphics[width=\textwidth]{rci_figure1.png}
\caption{%
\textbf{Left:} Cumulative KL divergence (AUC) per difficulty
tier; horizontal bars indicate group means.
\textbf{Centre:} Active expert coverage (\%) per task.
\textbf{Right:} Representative KL convergence curves;
dashed line indicates $\varepsilon{=}0.093$.}
\label{fig:results}
\end{figure}

\subsection{Model-Relative Complexity}

A notable finding is that MATH algebra problems (human label:
medium) require less cumulative KL divergence than GSM8K word
problems (human label: easy) for this model
($8.96$ vs.\ $10.73$, $p{<}0.001$).
We interpret this as evidence that RCI measures
\textbf{model-relative computational demand}---the amount of
information a specific model must extract from its own weight
structure to stabilise its output distribution on a specific
input---rather than task difficulty in the abstract.
This dissociation between model-internal and human-perceived
difficulty warrants further investigation across model families
and scales.

\subsection{Ablation}

Affinity-weighted BFS and random BFS yield equivalent average
steps on this pre-trained model without RCI-specific fine-tuning.
This is consistent with the hypothesis that the full benefit of
structured expansion requires training under the objective in
Equation~\ref{eq:loss}.

% ── 5. LIMITATIONS ────────────────────────────────────────
\section{Limitations and Future Work}
\label{sec:limitations}

\paragraph{Wall-clock latency.}
Each RCI step is a full forward pass; real-world speedup requires
fused sparse kernels.
Recent IO-aware GPU primitives~\citep{flashkmeans2025} suggest
this overhead may be substantially reducible without algorithmic
changes.

\paragraph{Scale.}
The PoC uses 1B active parameters.
Behaviour at 8B--70B scale remains unexamined.

\paragraph{Static affinity graph.}
Computed offline over 500 calibration examples; may be suboptimal
under distribution shift.
Dynamic affinity estimation is a natural extension.

\paragraph{Parallel heads.}
Multi-head parallelism requires multi-device infrastructure and
is not evaluated in the present PoC.

\paragraph{Convergence guarantees.}
RCI relies on empirical convergence.
Formal guarantees via contraction mapping are not pursued.

\paragraph{Future directions.}
(i) Attention-based pooling to mitigate oversmoothing at high
module coverage;
(ii) training from scratch under the full RCI objective;
(iii) validation across MoE families (Mixtral, Qwen-MoE,
DeepSeek-MoE);
(iv) extension to multi-modal architectures.

% ── 6. CONCLUSION ─────────────────────────────────────────
\section{Conclusion}
\label{sec:conclusion}

We introduced Recursive Convergent Inference (RCI), an inference
principle in which the set of active computational modules expands
monotonically from a minimal seed subset until empirical
convergence of the model's next-token output distribution.
In contrast to existing adaptive computation methods that operate
top-down---starting from a maximal graph and learning when to
halt---RCI operates bottom-up, growing computation as warranted
by the task rather than pruning it post-hoc.
Stopping requires no learned halting signal, external verifier,
or task-complexity pre-estimator; it emerges directly from the
distributional stability of the model's own output.

We instantiated RCI on OLMoE-1B-7B-0924, a fully open
mixture-of-experts model with 64 discrete experts, and provided
a reproducible implementation via Google Colab.
Empirical evaluation across reasoning benchmarks (GSM8K, MATH
algebra, MMLU hard subsets, $n{=}150$, 50 per tier) reveals two
findings.
First, all pairwise difficulty comparisons are statistically
significant: hard tasks accumulate significantly greater
cumulative KL divergence than easy ($U{=}1677$, $p{=}0.002$) and
medium tasks ($U{=}2106$, $p{<}0.001$), with a statistically
significant monotonic correlation overall
(Spearman $\rho{=}0.22$, $p{=}0.007$, $n{=}150$).
Second, and notably, RCI's complexity metric diverges from
human-defined difficulty labels: MATH algebra problems require
less cumulative KL divergence than GSM8K word problems for this
model ($8.96$ vs.\ $10.73$, $p{<}0.001$), suggesting an
intriguing dissociation between model-internal and
human-perceived difficulty that warrants further study across
architectures and scales.

We believe RCI offers a complementary perspective to the growing
literature on test-time compute scaling: that inference efficiency
can be achieved not only by generating fewer tokens or selecting
among multiple samples, but by dynamically allocating computation
at the level of parameter activation itself.

% ── REFERENCES ────────────────────────────────────────────
\bibliography{references}

\end{document}